\subsection{秩 1 紧群}

由《一般拓扑学》第 VIII 章 §1 第 4 节命题 3、命题 4 与注 4，拓扑群 \(\mathbf{SU}(2,\mathbf{C})\) 同构于范数为 1 的四元数组成的拓扑群 \(\mathbf{S}_3\) ，且 \(\mathbf{SU}(2,\mathbf{C})\) 关于由矩阵 \(I_2\) 与 \(-I_2\) 组成的子群 Z 的商同构于拓扑群 \(\mathbf{SO}(3,\mathbf{R})\) 。注意 Z 是 \(\mathbf{SU}(2,\mathbf{C})\) 的中心：事实上，由于 \(\mathbf{H} = \mathbf{R}.\mathbf{S}_3\) ，群 \(\mathbf{S}_3\) 的中心的每个元素属于中心 \(\mathbf{R}\) （代数 \(\mathbf{H}\) 的） ，从而属于二元素群 \(\mathbf{S}_3 \cap \mathbf{R} = \{-1,1\}\) 。

命题 6. 秩 1 的每个紧半单实李代数同构于 \(\mathfrak{su}(2,\mathbf{C})\) 与 \(\mathfrak{o}(3,\mathbf{R})\) 。秩 1 的每个连通半单紧李群，若单连通则同构于 \(\mathbf{SU}(2,\mathbf{C})\) ，否则同构于 \(\mathbf{SO}(3,\mathbf{R})\) 。

第一个断言由命题 4 的推论与命题 5 得出。由于 \(\mathbf{SU}(2,\mathbf{C})\) 同胚于 \(\mathbf{S}_3\) （《一般拓扑学》，第 VIII 章，§1，第 4 节，注 4），从而单连通（《一般拓扑学》，第 XI 章，编写中），故每个单连通秩 1 紧半单李群同构于 \(\mathbf{SU}(2,\mathbf{C})\) ；每个不单连通的秩 1 连通紧半单李群同构于 \(\mathbf{SU}(2,\mathbf{C})\) 关于子群 \(Z\) （不约化为单位元的）的商，即同构于 \(\mathbf{SO}(3,\mathbf{R})\) 。

注. 我们在上面已经看到 \(\mathbf{SU}(2,\mathbf{C})\) 单连通，而 \(\pi_{1}(\mathbf{SO}(3,\mathbf{R}))\) 的阶为 2。我们将在很后面看到这些结果分别推广到 \(\mathbf{SU}(n,\mathbf{C})\) （\(n\geq1\) ）与 \(\mathbf{SO}(n,\mathbf{R})\) （\(n\geq3\) ）（另见 §3 习题 4 与 5）。

回忆（第 VIII 章，§1，第 1 节）：\(\mathfrak{sl}(2,\mathbf{C})\) 的典则基是基 \((X_{+},X_{-},H)\) ，其中

\[
X _ {+} = \left( \begin{array}{c c} 0 & 1 \\ 0 & 0 \end{array} \right), X _ {-} = \left( \begin{array}{c c} 0 & 0 \\ - 1 & 0 \end{array} \right), H = \left( \begin{array}{c c} 1 & 0 \\ 0 & - 1 \end{array} \right).
\]

我们由此得到一个基 \((U, V, iH)\) （\(\mathfrak{su}(2, \mathbf{C})\) 的），也称它为典则基，其定义是

\[
\begin{array}{c} U = X _ {+} + X _ {-} = \left( \begin{array}{c c} 0 & 1 \\ - 1 & 0 \end{array} \right), V = i (X _ {+} - X _ {-}) = \left( \begin{array}{c c} 0 & i \\ i & 0 \end{array} \right), \\ i H = \left( \begin{array}{c c} i & 0 \\ 0 & - i \end{array} \right). \end{array}
\]

我们有

\[
[ i H, U ] = 2 V, [ i H, V ] = - 2 U, [ U, V ] = 2 i H.\tag{13}
\]

若 B 表示 \(\mathfrak{su}(2,\mathbf{C})\) 的 Killing 型，则直接计算给出

\[
\mathrm{B} (a U + b V + c i H, a ^ {\prime} U + b ^ {\prime} V + c ^ {\prime} i H) = - 8 (a a ^ {\prime} + b b ^ {\prime} + c c ^ {\prime}),\tag{14}
\]

于是，若借助典则基把 \(\mathfrak{su}(2,\mathbf{C})\) 与 \(R^{3}\) 等同起来，则 \(\mathbf{SU}(2,\mathbf{C})\) 的伴随表示定义一个同态 \(\mathbf{SU}(2,\mathbf{C})\to\mathbf{SO}(3,\mathbf{R})\) （参见上文）。

进而注意，\(\mathbf{Ri}H\) 是 \(\mathfrak{su}(2,\mathbf{C})\) 的嘉当子代数，与之对应的极大环面 \(T\) （\(\mathbf{SU}(2,\mathbf{C})\) 的）由对角矩阵 \(\left( \begin{array}{cc}a & 0\\ 0 & \bar{a} \end{array} \right)\) （其中 \(a\bar{a} = 1\) ）组成，且指数映射

\(\exp : \mathbf{R}iH \to \mathrm{T}\)

把 \(xH\) （\(x \in \mathbf{R}i\) ）映为矩阵 \(\left( \begin{array}{cc}\exp (x) & 0\\ 0 & \exp (-x) \end{array} \right)\) ，故其核为 \(\mathbf{Z}.K\) ，其中 \(K\) 是 \(\mathfrak{su}(2,\mathbf{C})\) 中由下式定义的元素

\[
K = 2 \pi i H = \left( \begin{array}{c c} 2 \pi i & 0 \\ 0 & - 2 \pi i \end{array} \right).\tag{15}
\]

进而，\(\mathbf{SU}(2,\mathbf{C})\) 的中心由单位元与 \(\exp (K / 2)\) 组成。

\[
\theta = \left( \begin{array}{c c} 0 & - 1 \\ 1 & 0 \end{array} \right) \in \mathbf {S U} (2, \mathbf {C}).\tag{16}
\]

由第 VIII 章 §1 第 5 节，

\[
\theta^ {2} = \left( \begin{array}{c c} - 1 & 0 \\ 0 & - 1 \end{array} \right), (\text {Int}   \theta) t = t ^ {- 1}, t \in \text {T},\tag{17}
\]

\[
(\operatorname{Ad} \theta) X _ {+} = X _ {-}, (\operatorname{Ad} \theta) X _ {-} = X _ {+}, (\operatorname{Ad} \theta) U = U, (\operatorname{Ad} \theta) V = - V.\tag{18}
\]

\[
\text {   Finally,   for   } t = \left( \begin{array}{c c} a & 0 \\ 0 & \bar {a} \end{array} \right) \in \mathrm{T}, \text {   we   have   }
\]

\[
(\operatorname{Ad} t) X _ {+} = a ^ {2} X _ {+}, (\operatorname{Ad} t) X _ {-} = a ^ {- 2} X _ {-}, (\operatorname{Ad} t) H = H,\tag{19}
\]

\[
(\operatorname{Ad} t) U = \mathcal {R} (a ^ {2}) U + \mathcal {I} (a ^ {2}) V, (\operatorname{Ad} t) V = - \mathcal {I} (a ^ {2}) U + \mathcal {R} (a ^ {2}) V.\tag{20}
\]

